	\section{Kostant--Whittaker reduction and quantization}
	\label{sect::kw_red}
	
	In this section, we consider the Whittaker analog of the parabolic reduction functor from Section~\ref{subsect::parabolic_reduction} that we call (following \cite{BezrukavnikovFinkelberg}) the \emph{Kostant--Whittaker reduction}. This is a finite-dimensional analog of the Drinfeld--Sokolov reduction, see \cite{Arakawa}.
	
	\subsection{Mirabolic setting} Consider the category $\HC(\rM_N)$ of Harish-Chandra bimodules over the mirabolic subalgebra~$\m_N$ as in Section~\ref{subsect::hc_bimod}. Denote by $\Wh(\m_N)$ the category of right Whittaker modules over $\m_N$ from Definition~\ref{def::whittaker_module}. It has a distinguished object \smash{$Q_{\m_N} = \n_-^{\psi} \backslash \U(\m_N)$}. Similarly to \eqref{eq::act_h}, we can define a functor
	\begin{equation*}
		\act_{\C[\hbar]}^{\psi} \colon\ \Mod{\C[\hbar]} \rightarrow \Wh(\m_N), \qquad A \mapsto A \otimes_{\C[\hbar]} Q_{\m_N}.
	\end{equation*}
	Then we have an analog of Kostant's equivalence \cite{Kostant}:
	\begin{Theorem}[mirabolic Kostant's equivalence]
		\label{thm::mirabolic_skryabin}
		The functor $\act_{\C[\hbar]}^{\psi}$ is an equivalence.
		\begin{proof}
			We proceed as in Proposition~\ref{prop::act_equiv}. The functor $(-)^{N_-} \colon \Wh(\m_N) \rightarrow \Mod{\C[\hbar]}$ is right adjoint to \smash{$\act_{\C[\hbar]}^{\psi}$}
			\begin{align*}
				\Hom_{\Wh(\m_N)} \bigl(\act_{\C[\hbar]}^{\psi}(A), M\bigr) &= \Hom_{\Wh(\m_N)}(A \otimes_{\C[\hbar]} Q_{\m_N},M) \\
				&\cong \Hom_{\U(\m_N)} (A \otimes_{\C[\hbar]} Q_{\m_N},M) \\
				&\cong \Hom_{\C[\hbar]}\bigl(A,M^{N_-}\bigr).
			\end{align*}
			The unit of the adjunction is given by
$
			A \mapsto (A \otimes_{\C[\hbar]} Q_{\m_N})^{N_-}$.
			The Kirillov projector gives an isomorphism with the quotient functor as in \eqref{eq::kirillov_proj_iso_inv_quot}
			\[
			(-)_{\b^{\vec{u}}}\colon\ \Wh(\m_N) \rightarrow \Mod{\C[\hbar]}, \qquad M \mapsto M/\b^{\vec{u}}.
			\]
			So, the functor $(-)^{N_-}$ is exact. We have
			\[
			(A \otimes_{\C[\hbar]} Q_{\m_N})^{N_-} \cong (A \otimes_{\C[\hbar]} Q_{\m_N})/ \b^{\vec{u}}
			\]
			and the composition
			\[
			A \mapsto (A \otimes_{\C[\hbar]} Q_{\m_N})^{N_-} \cong (A \otimes_{\C[\hbar]} Q_{\m_N})/ \b^{\vec{u}}
			\]
			is an isomorphism by the PBW theorem. The rest follows by the same arguments as in Proposition~\ref{prop::act_equiv}.
		\end{proof}
	\end{Theorem}
	
	Similarly to \eqref{eq::act_g}, one can define the action functor
	\[
	\act_{\m_N}^{\psi}\colon\ \HC(\rM_N) \rightarrow \Wh(\m_N), \qquad X \mapsto Q \otimes_{\U(\m_N)} X.
	\]
	The following definition is an analog of Definition~\ref{def::parabolic_reduction}:
	\begin{Definition}
		The \emph{mirabolic Kostant--Whittaker reduction} functor $\res_{\m_N}^{\psi}  \colon  \HC(\rM_N)  \rightarrow \ \allowbreak\Mod{\C[\hbar]}$ is the composition
		\[
		\HC(\rM_N) \xrightarrow{\act_{\m_N}} \Wh(\m_N) \xrightarrow{(-)^{N_-}} \Mod{\C[\hbar]}.
		\]
	\end{Definition}
	
	Explicitly, it is given by a quantum Hamiltonian reduction \smash{$X \mapsto \bigl(\n_-^{\psi} \backslash X\bigr)^{N_-}$}. There is a~natural lax monoidal structure on \smash{$\res_{\m_N}^{\psi}$}
	\[
	\bigl(\n_-^{\psi} \backslash X\bigr)^{N_-} \otimes_{\C[\hbar]} \bigl(\n_-^{\psi} \backslash Y\bigr)^{N_-} \rightarrow \bigl(\n_-^{\psi} \backslash X\otimes_{\U(\m_N)} Y\bigr)^{N_-}.
	\]
	As in Theorem~\ref{thm::res_exact_monoidal}, we have the following.
	\begin{Theorem}
		The functor $\res_{\m_N}^{\psi}$ is colimit-preserving and monoidal.
		\begin{proof}
			Follows from exactly the same arguments as \cite[Corollary 4.18]{KalmykovSafronov}.
		\end{proof}
	\end{Theorem}
	
	Recall the monoidal functor $\free \colon \Rep(\rM_N) \rightarrow \HC(\rM_N)$ from Section~\ref{subsect::hc_bimod}. On the one hand, we get a monoidal functor $\Rep(\rM_N) \rightarrow \Mod{\C[\hbar]}$ by composing with \smash{$\res_{\m_N}^{\psi}$}. On the other hand, there is the composition of the forgetful functor with the extensions of scalars from $\C$ to $\C[\hbar]$ that we denote by $\forget\colon \Rep(\rM_N) \rightarrow \Mod{\C[\hbar]}$. So, we have a diagram
	\begin{equation}
		\label{eq::mirabolic_res_diag}
\begin{split}
& \xymatrix{
			\Rep(\rM_N) \ar[r]^{\free} \ar[dr]_{\forget} & \HC(\rM_N) \ar[d]^{\res_{\m_N}^{\psi}} \\
			& \Mod{\C[\hbar]}.
		}
\end{split}
	\end{equation}
	The following result is an analog of Theorem~\ref{thm::dynamical_trivialization}.
	\begin{Theorem}
		\label{thm::mirabolic_kw_trivialization}
		The Kirillov projector defines a natural isomorphism
		\[
		P_{\m_N}(\vec{u})_V \colon\ V \rightarrow \bigl(\n_-^{\psi} \backslash V \otimes \U(\m_N)\bigr)^{N_-}.
		\]
		In other words, the diagram \eqref{eq::mirabolic_res_diag} is commutative.
		\begin{proof}
			By \eqref{eq::kirillov_proj_iso_inv_quot}, we have
			\[
			\bigl(\n_-^{\psi} \backslash V \otimes \U(\m_N)\bigr)^{N_-} \cong \n_-^{\psi} \backslash V \otimes \U(\m_N) / \b^{\vec{u}}.
			\]
			By the PBW theorem, the composition
			\[
			V \rightarrow \bigl(\n_-^{\psi} \backslash \U(\m_N) \otimes V\bigr)^{N_-} \cong \n_-^{\psi} \backslash V \otimes \U(\m_N) / \b^{\vec{u}}
			\]
			is an isomorphism.
		\end{proof}
	\end{Theorem}
	
	
	In particular, there is a monoidal structure on $\forget$ induced from $\res_{\m_N}^{\psi}$.
	\begin{Theorem}\quad
		\label{thm::mirabolic_tensor_structure}
		\begin{enumerate}\itemsep=0pt
			\item[$(1)$] There is a collection of maps $F_{VW}(\vec{u})\in \End_{\C[\hbar]}(V[\hbar] \otimes_{\C[\hbar]} W[\hbar])$ natural in $V,W\in \Rep(\rM_N)$, such that the diagram
			\[
			\xymatrixcolsep{5pc}\xymatrix{
				V[\hbar] \otimes_{\C[\hbar]} W[\hbar] \ar[r]^{F_{VW}(\vec{u})} \ar[d]_{(P_{\m_N}(\vec{u}))_V \otimes (P_{\m_N}(\vec{u}))_W} & V[\hbar] \otimes_{\C[\hbar]} W[\hbar] \ar[d]^{(P_{\m_N}(\vec{u}))_{V\otimes W}} \\
				V[\hbar] \otimes_{\C[\hbar]} W[\hbar] \ar[r] & V[\hbar] \otimes_{\C[\hbar]} W[\hbar]
			}
			\]
			is commutative, where the lower right arrow is the natural tensor structure on $\res_{\m_N}^{\psi}$. In particular, the collection of $F_{VW}(\vec{u})$ satisfies the twist equation \eqref{eq::twist}, and $R^{\CG}_{VW} (\vec{u}):= F_{WV}(\vec{u})^{-1} F_{VW}(\vec{u})$ is a solution to the quantum Yang--Baxter equation.
			\item[$(2)$] The map $F_{VW}(\vec{u})$ has the form
			\[
			F_{VW} (\vec{u})\in \id_{V\otimes W} + \hbar \mathrm{U}(\n_-)^{>0} \otimes \mathrm{U}(\b)^{>0},
			\]
			where the upper subscript $>0$ means the augmentation ideal.
			\item[$(3)$] The inverse $F_{VW}(\vec{u})^{-1}$ also satisfies
			\[
			F_{VW}(\vec{u})^{-1} \in \id_{V\otimes W} + \hbar \mathrm{U}(\n_-)^{>0} \otimes \mathrm{U}(\b)^{>0}.
			\]
		\end{enumerate}
		\begin{proof}
			We proceed exactly as in Theorem~\ref{thm::dynamical_quantization}. For brevity, denote $P = P_{\m_N}(\vec{u})$. Every element in \smash{$\bigl(\n_-^{\psi} \backslash V\otimes \U(\m_N) \bigr)^{N_-}$} can be presented as $PvP$ for $v\in V$ by Theorem~\ref{thm::mirabolic_kw_trivialization}. Therefore, we need to show that
			\begin{equation}
				\label{eq::middle_p_twist}
				PvPwP = PF(v\otimes w) P.
			\end{equation}
			
			First, let us choose the following PBW basis: $\bigl\{E_{ij}^{\psi}\bigr\}$ for $\U(\n_-)$ and
			\[
			(E_{kl}, E_{kk} + u_k \mid 1 \leq k < l \leq N-1)
			\]
			for $\U(\b)$ (exactly in this order for the latter). It is clear from the construction \eqref{eq::kirillov_proj} and Remark~\ref{rmk::kirillov_proj_image_quot} that
			\begin{equation}
				\label{eq::kirillov_proj_normalization}
				P \in 1 + \n_-^{\psi} \U(\m_N) \cap \U(\m_N) \b^{\vec{u}}
			\end{equation}
			(recall the notation $\b^{\vec{u}}$ from \eqref{eq::borel_u}). Let us write it in this PBW basis
			\begin{equation*}
				P = 1 + \hbar^{k_i}\sum\limits_i f^{\psi}_i e_i^u,
			\end{equation*}
			where $f^{\psi}_i \in \U(\n_-)$, $e_i^u \in \U(\b)$, and $k_i$ is some negative number. This is well defined, since a~part of $P$ that acts non-trivially is actually finite. It is clear from \eqref{eq::kirillov_proj_normalization} that $e^u_i=1$ if and only if $f^{\psi}_i=1$.
			
			Now consider the middle $P$ in \eqref{eq::middle_p_twist}. As in Theorem~\ref{thm::dynamical_quantization}, we push the $f_i^{\psi}$-term to the left until it meets $P$ and becomes zero. Likewise, we push the $e_i^u$-term to right until it meets $P$ and becomes zero as well. Since every term
			\[
			\hbar^{-k} \bigl(-E_{ij}^{\psi}\bigr) (-1)^{(i+j)k} \frac{\prod_{i=0}^{k-1} L_{j \dots i-1}^{j \dots i-1}(u-\hbar i + \hbar)}{k!}
			\]
			in \eqref{eq::kirillov_proj_term} acts with the power at least $\hbar^k$, the second part of the theorem follows. The third part follows from the fact that $F_{VW}(\vec{u})$ is upper-triangular in the natural filtration on $V\otimes W$ induced from a $\n_-$-filtration on $V$ and $\b$-filtration on $W$.
		\end{proof}
	\end{Theorem}
	
	\begin{Remark}
		\label{rmk::mirabolic_tensor_structure}
		In general, let $X$ be any Whittaker module and $V$ any representation of $\rM_N$. Then for any $x\in X$ and $v\in V$,
		\[
		PxvP = \sum\limits_i Px_i P v_i P
		\]
		for some $x_i\in X$, $v_i \in V$. Indeed, recall that by Definition~\ref{def::whittaker_module}, the action of $\hbar^{-1} \n_-^{\psi}$ integrates to that of $N_-$. Therefore, we can proceed exactly as in the proof of Theorem~\ref{thm::mirabolic_tensor_structure} and use the fact that the analog of $F_{VW}(\vec{u})$ in this case is invertible as well.
	\end{Remark}
